\textbf{1. Kiến thức cần nhớ}
\begin{itemize}
\item \textbf{Định lí Viète}: Nếu phương trình $ax^2 + bx + c = 0$ ($a \ne 0$) có hai nghiệm $x_1, x_2$ thì:
$$\boxed{\begin{cases} S = x_1 + x_2 = -\dfrac{b}{a} \\ P = x_1 x_2 = \dfrac{c}{a} \end{cases}}$$
\item \textbf{Trường hợp nhẩm nghiệm đặc biệt}:
\begin{itemize}
\item Nếu $a + b + c = 0$ thì phương trình có nghiệm $\boxed{x_1 = 1;\quad x_2 = \dfrac{c}{a}}$.
\item Nếu $a - b + c = 0$ thì phương trình có nghiệm $\boxed{x_1 = -1;\quad x_2 = -\dfrac{c}{a}}$.
\end{itemize}
\item \textbf{Tìm hai số biết tổng và tích}: Nếu hai số có tổng bằng $S$ và tích bằng $P$ thì hai số đó là nghiệm của phương trình:
$$\boxed{X^2 - SX + P = 0}\quad (\text{điều kiện tồn tại: } S^2 - 4P \ge 0)$$
\end{itemize}

\textbf{2. Phương pháp giải các dạng}

%% DANG 9D6?20-1
\textbf{Dạng 1. Tính tổng, tích hai nghiệm; nhẩm nghiệm}
\begin{enumerate}
\item Kiểm tra phương trình có nghiệm ($\Delta \ge 0$ hoặc $a \cdot c < 0$).
\item Áp dụng trực tiếp: $x_1 + x_2 = -\dfrac{b}{a}$ và $x_1 x_2 = \dfrac{c}{a}$.
\item Kiểm tra nhẩm nghiệm: xét $a + b + c$ hoặc $a - b + c$.
\end{enumerate}

%% DANG 9D6?20-2
\textbf{Dạng 2. Tính giá trị biểu thức đối xứng của hai nghiệm}
\begin{enumerate}
\item Các biến đổi đối xứng kinh điển sang $S$ và $P$:
\begin{itemize}
\item $x_1^2 + x_2^2 = (x_1 + x_2)^2 - 2x_1 x_2 = \boxed{S^2 - 2P}$.
\item $(x_1 - x_2)^2 = (x_1 + x_2)^2 - 4x_1 x_2 = \boxed{S^2 - 4P}$.
\item $x_1^3 + x_2^3 = (x_1 + x_2)^3 - 3x_1 x_2(x_1 + x_2) = \boxed{S^3 - 3SP}$.
\item $\dfrac{1}{x_1} + \dfrac{1}{x_2} = \dfrac{x_1 + x_2}{x_1 x_2} = \boxed{\dfrac{S}{P}}$.
\item $\dfrac{x_1}{x_2} + \dfrac{x_2}{x_1} = \dfrac{x_1^2 + x_2^2}{x_1 x_2} = \boxed{\dfrac{S^2 - 2P}{P}}$.
\end{itemize}
\item Thay $S = -\dfrac{b}{a}$ và $P = \dfrac{c}{a}$ vào để tính kết quả số.
\end{enumerate}

%% DANG 9D6?20-3
\textbf{Dạng 3. Tìm hai số biết tổng và tích}
\begin{enumerate}
\item Cho biết tổng $S$ và tích $P$. Kiểm tra điều kiện có hai số: $S^2 - 4P \ge 0$.
\item Lập phương trình bậc hai $X^2 - SX + P = 0$.
\item Giải phương trình tìm hai nghiệm $X_1, X_2$, từ đó kết luận hai cặp số cần tìm.
\end{enumerate}

%% DANG 9D6?20-4
\textbf{Dạng 4. Tìm tham số thỏa mãn điều kiện về nghiệm}
\begin{enumerate}
\item \textbf{Bước 1}: Tìm điều kiện để phương trình có hai nghiệm: $\boxed{\begin{cases} a \ne 0 \\ \Delta \ge 0 \end{cases}}$ (hoặc $\Delta > 0$ nếu đề yêu cầu hai nghiệm phân biệt).
\item \textbf{Bước 2}: Viết hệ thức Vi-ét theo tham số $m$: $S(m)$ và $P(m)$.
\item \textbf{Bước 3}: Biến đổi biểu thức điều kiện của đề bài theo $S$ và $P$, giải phương trình tìm $m$.
\item \textbf{Bước 4}: \textbf{Đối chiếu với điều kiện ở Bước 1} rồi kết luận $m$.
\end{enumerate}
